\subsection{CX-N15: a complete accessible composition for general factorization}
\label{finding:CX-N15}

\paragraph{Finding and source interface.}
Assume \(\mathsf H\), the all-number-field Hecke claim.
Given an explicit finite field \(\mathbb F_q\), \(q=p^m\), and
\(0\ne f\in\mathbb F_q[X]\) of degree \(d\geq2\), there is a
deterministic algorithm that returns its complete factorization in
\[
 \bigl(d^{\log d}\log q\bigr)^{O(1)}
\]
bit operations.
The composition below uses Adleman--Lenstra Algorithm~A and
Theorem~1, with its auxiliary-prime step supplied by CX-N11;
Bach--von zur Gathen--Lenstra Theorem~2; and the IKRS Main Theorem,
together with the explicitly proved coefficient-unit step of
CX-N14
\cite{CXNAdlemanLenstra86,CXNBGL01,CXNIKRS08}.
It does not require the unretrieved internal reduction of
Evdokimov's 1994 paper, or a general finite-field factorization
oracle.
This gives an independent proof of the broad transfer in CX-N09;
it does not change that theorem's degree dependence.

\paragraph{Step 1: construct a maximal coefficient unit.}
Consider the unresolved prime-order branch of CX-N14, with
\(p>d_0\geq r\), where \(d_0\) is the dimension of its split
algebra and \(r\) is prime. Put
\[
 s=\operatorname{ord}_r(p),\qquad Q=p^s,
 \qquad e=v_r(Q-1),\qquad R=r^e.
\]
The quantities \(s\leq r-1\), \(e\), and \(R\) are computable
in polynomial time in \(r,\log p\); no factorization of the
complementary part of \(Q-1\) is used.
Adleman--Lenstra Algorithm~A, with CX-N11, constructs an explicit
field
\(U=\mathbb F_{p^{sr}}\) in polynomial time in \(r,\log p\).
Indeed \(sr\leq r(r-1)\) and \(p\nmid sr\), because both
\(r\) and \(s\) are smaller than \(p\).
Compute the subfield
\[
 B=\ker\bigl(U\longrightarrow U,\ u\longmapsto u^Q-u\bigr)
   \cong\mathbb F_Q
\]
by \(\mathbb F_p\)-linear algebra, and compute its multiplication
table and inclusion in \(U\).
This is exactly the explicit-data convention of BGL, adopted
from Lenstra: basis structure constants suffice, and Frobenius
kernels provide explicit subfields
\cite[p.~329 and \S2, p.~332]{CXNLenstra91}.

BGL Theorem~2, p.~7, constructs a primitive \(r\)-th root
\(\zeta\in B\), given additionally an irreducible polynomial
over \(\mathbb F_p\) of degree \(\ell\) for each prime
\(\ell\mid r-1\) with \(\ell\nmid s\).
All these primes are smaller than \(p\) and at most \(r-1\).
Trial division factors \(r-1\) in polynomial time in \(r\),
and Adleman--Lenstra with CX-N11 constructs every required
auxiliary polynomial in polynomial time in \(r,\log p\).
Thus all BGL inputs are supplied, and its output and running-time
conditions apply.

Find a nonzero solution of
\[
 u^Q=\zeta u\qquad (u\in U)
\]
by \(\mathbb F_p\)-linear algebra.
Such a solution exists: the image of the power map
\(U^\times\to U^\times\), \(v\mapsto v^{Q-1}\), is the
cyclic subgroup of order
\[
 S=\frac{Q^r-1}{Q-1}=1+Q+\cdots+Q^{r-1}.
\]
Since \(Q\equiv1\pmod r\), the integer \(S\) is divisible
by \(r\); hence this subgroup contains \(\zeta\).
Let \(c=u^r\). Then
\[
 c^Q=(\zeta u)^r=c,
 \qquad c^{(Q-1)/r}=u^{Q-1}=\zeta\ne1.
\]
Thus \(c\in B\) is an \(r\)-th nonresidue.
Because \(B^\times\) is cyclic of order \(Q-1=R w_0\),
\(r\nmid w_0\), the element
\[
 \eta_0=c^{w_0}\in B^\times
\]
has exact order \(R\).
To verify the last assertion, write \(c=g^a\) for an abstract
generator \(g\) of \(B^\times\).
The nonresidue test says \(r\nmid a\), so
\(\operatorname{ord}(g^{a w_0})=R\).
This argument uses a generator only to prove the order; none is
an algorithmic input.

\paragraph{Step 2: place the unit in the cyclotomic coefficient algebra.}
Let
\[
 D_r=\mathbb F_p[T]/(\Phi_r(T)),\qquad
 \rho_b(T)=T^b\quad\bigl(b\in G=(\mathbb Z/r\mathbb Z)^\times\bigr).
\]
Enumerate the powers \(\zeta^a\), \(a\in G\), and their
Frobenius orbits under \(a\mapsto pa\pmod r\).
Every orbit \(O\) has size \(s\), and
\[
 g_O(T)=\prod_{a\in O}(T-\zeta^a)
\]
has coefficients in \(\mathbb F_p\) and is irreducible of
degree \(s\).
These statements follow because \(O\) is precisely the
Frobenius orbit of any one of its roots.
Computing the products gives the full factorization
\(\Phi_r=\prod_Og_O\) explicitly; no polynomial-factorization
oracle has been used.

Let \(O_0\) be the orbit containing \(1\).
The evaluation map
\(\mathbb F_p[T]/(g_{O_0})\to B\), \(T\mapsto\zeta\), is
an isomorphism: \(\zeta\) has degree \(s\) and generates \(B\).
Its inverse is computed by expressing \(\eta_0\) on the basis
\(1,\zeta,\ldots,\zeta^{s-1}\).
The Chinese remainder theorem, implemented with polynomial
Euclidean algorithms, now constructs \(w\in D_r^\times\)
whose value on this factor is \(\eta_0\) and whose value on
all other factors is \(1\).
In particular \(w^R=1\).

\paragraph{Step 3: a projector with a proved order bound.}
For \(b\in G\), put
\[
 \omega(b)=b^{r^{e-1}}\bmod R,
 \qquad \kappa=(r-1)^{-1}\bmod R,
 \qquad \nu_b=\kappa\,\omega(b)^{-1}\bmod R.
\]
Here \(\omega\) is the multiplicative Teichm\"uller lift of
\(G\): it reduces to \(b\) modulo \(r\), and its values have
order dividing \(r-1\).
Define, using exponents reduced modulo \(R\),
\[
 \eta=\prod_{b\in G}\rho_b(w)^{\nu_b}\in D_r^\times.
\]
All factors have \(r\)-power order dividing \(R\), and the
product is computable in polynomial time.
For \(a\in G\), reindexing by \(b'=ab\) gives
\[
 \rho_a(\eta)
 =\prod_{b'\in G}\rho_{b'}(w)^{{\kappa}\omega(a)\omega(b')^{-1}}
 =\eta^{\omega(a)}.
\]
Consequently \(\eta\) lies in the coefficient Teichm\"uller
subgroup required by CX-N14.

It remains to prove that projection has not decreased its order.
Evaluate at \(T=\zeta\), the factor indexed by \(O_0\).
Only the terms \(b\in\langle p\rangle\leq G\) contribute
nonidentity values. Write them as \(b=p^j\pmod r\),
\(0\leq j<s\). Since \(w\) has coefficients in
\(\mathbb F_p\),
\[
 \rho_{p^j}(w)(\zeta)=w(\zeta^{p^j})=\eta_0^{p^j}.
\]
The congruence \(p^s\equiv1\pmod R\), together with
\(s\mid r-1\), implies that \(p\bmod R\) is the
Teichm\"uller lift of \(p\bmod r\). Indeed reduction modulo
\(r\) is injective on the prime-to-\(r\) torsion subgroup of
\((\mathbb Z/R\mathbb Z)^\times\), and \(p\) has order
\(s\) in that subgroup.
Hence \(\omega(p^j)=p^j\pmod R\), and the evaluation is
\[
 \eta(\zeta)=\eta_0^{\kappa s}.
\]
Both \(s\) and \(r-1\) are coprime to \(r\), so
\(\kappa s\) is coprime to \(r\).
This component has order \(R\), proving
\(\operatorname{ord}(\eta)=R\).
The required additional input of CX-N14 is therefore constructed
under \(\mathsf H\), in \(\operatorname{poly}(d_0,\log p)\)
time. Together with the other cases of CX-N14, this splits any
nontrivial split prime-field algebra for which a nonidentity
automorphism is supplied.

\paragraph{Step 4: from a general polynomial to a split algebra.}
Polynomial gcds and derivatives give squarefree preprocessing.
If a positive-degree polynomial has zero derivative, it is a
\(p\)-th power: take the coefficient \(p\)-th roots using the
inverse Frobenius \(a\mapsto a^{p^{m-1}}\), divide the exponents
by \(p\), and retain the multiplicity \(p\).
Otherwise a nontrivial gcd with the derivative yields smaller
factors to process recursively.
These operations have polynomial bit complexity.
It therefore suffices to factor a monic squarefree polynomial
\(f\) of degree at most \(d\).

Compute the multiplication table of
\(A=\mathbb F_q[X]/(f)\), viewed as an
\(\mathbb F_p\)-algebra of dimension \(m\deg f\), and then
compute
\[
 B_f=\ker\bigl(A\longrightarrow A,\ a\longmapsto a^p-a\bigr).
\]
The map is \(\mathbb F_p\)-linear, so its kernel and the kernel's
multiplication table are obtained by linear algebra.
Abstractly, if \(f=\prod_{i=1}^t f_i\) is its distinct
irreducible factorization, then
\[
 A\cong\prod_{i=1}^t\mathbb F_{q^{\deg f_i}},
 \qquad B_f\cong\mathbb F_p^t.
\]
This proves the split promise without knowing the factors.
The primitive idempotents of \(B_f\) and \(A\) coincide.
If \(\dim_{\mathbb F_p}B_f=1\), the polynomial is irreducible.

Otherwise apply the IKRS Main Theorem, p.~3, to the explicit
\(t\)-dimensional algebra \(B_f\).
It deterministically gives a decomposition into ideals, each
with an automorphism whose order equals that ideal's dimension,
in time \(\operatorname{poly}(t^{\log t},\log p)\).
If the decomposition has more than one summand, the identity of
any proper nonzero summand is a nonzero zero divisor of \(B_f\);
that identity is found by linear algebra in the supplied ideal.
If there is only one summand, its automorphism has order
\(t>1\), and the preceding steps with CX-N14 produce a nonzero
zero divisor of \(B_f\) in polynomial additional time.
Thus in either case we obtain a nonzero zero divisor \(b\in A\),
and \(\gcd(f,b)\), using its polynomial representative, is a
proper nonconstant factor.
Equivalently, a primitive idempotent \(e_i\) identifies the
component factor as \(\gcd(f,1-e_i)\).
Recurse until all factors are irreducible, and reinsert the
recorded multiplicities.

\paragraph{Running time and exact analytic scope.}
There are at most \(2d-1\) nodes in the nontrivial factorization
tree, and every split algebra has dimension at most \(d\).
Linear algebra in the original quotient has size polynomial in
\(d,m,\log p\), hence in \(d,\log q\).
At each node the coefficient construction has polynomial cost;
the IKRS call has cost
\(d^{O(\log d)}(\log q)^{O(1)}\).
Multiplying by the number of nodes preserves the displayed
\((d^{\log d}\log q)^{O(1)}\) bound.

The analytic assumptions enter the constructed fields and
auxiliary irreducible polynomials through Adleman--Lenstra's
auxiliary-prime step only. CX-N11 supplies that step from its
explicit Dedekind-zeta hypotheses, all implied by \(\mathsf H\).
For this composition the required degrees are \(sr\) and the
specified primes \(\ell\mid r-1\); they are bounded by \(d^2\)
and are coprime to \(p\) in the branch where they are used.
BGL's root-of-unity algorithm with those inputs, the
Frobenius/eigenvector calculations, CRT, the projector, and IKRS
are unconditional algebraic steps.
The construction therefore supplies the formerly isolated
coefficient-unit gate under \(\mathsf H\) without circularly
calling general factorization or assuming a polynomial bound
for its varying degree.
\hfill\(\square\)

\paragraph{Finite-\'etale restatement.}
For squarefree \(f\), \(\operatorname{Spec}A\) is a finite
\'etale scheme over \(\mathbb F_p\).
Its Frobenius-fixed function algebra \(B_f\) is the constant
finite \'etale algebra indexing its connected components.
Idempotents found in \(B_f\) pull back to the corresponding
open-and-closed decomposition of \(\operatorname{Spec}A\).
The proof constructs these component projectors with the stated
quasipolynomial dependence on the input polynomial degree.
